\begin{abstract}
Integration-by-parts reduction expresses Feynman integrals in a master basis~\cite{KiraOriginal}. We combine learned seed selection with exact preprocessing and compare their separate effects on complete reduction time. Validation selects a 48-tree extremely randomized forest with maximum depth 10 from eleven fitted cost models and a preceding neural model. Its 44 structural inputs predict the costs of ten seed settings. The corrected test defines 92 requests, and the learned application completes three exact checks on 91. For 63 matched requests that invoke ML, the summed median times are 208.698 seconds with the model, 210.805 with native choices, and 210.587 with the tuned fixed choice. Both repetition intervals include equal time. Including the unfinished request, total attempted time is higher with ML. The derived vacuum formula gives 5.35 times faster reductions on 4 matched requests than the initial block interface; the massless formula gives 1.28 times faster reduction on 1 matched request. These analytic gains precede ML. Exact factorization and the derived formulas extend specified families through thirteen loops. The products contain small independent blocks, and the formulas cover one-mass vacuum and massless two-point integrals with one quadratic numerator. General irreducible high-loop scattering remains unverified. Timings include prediction, equation generation, reduction, checks, and recovery, with basis preparation recorded separately. Native comparisons keep the dimension symbolic and otherwise use rational kinematics. Earlier local seed and equation tests show measured savings; learned elimination-row selection remains omitted because it gives no reliable complete-time gain.
\end{abstract}
